%------------------------------------------------------------------------------%
\documentclass[fleqn,utf8,aspectratio=169,12pt]{beamer}

\usepackage{gaal}

\title{Bases Ortogonias e Ortonormais}

%------------------------------------------------------------------------------%
\begin{document}

\addtitle

%------------------------------------------------------------------------------%
\section{Bases e Coordenadas}

%------------------------------------------------------------------------------%
\begin{frame}{Base}
  Uma base de um espaço vetorial $V$
  \[
    \pause
    \mathcal{B} = \{\vec{v}_1,\ldots,\vec{v}_n\}
  \]
  \pause
  é um conjunto de vetores que
  \pause
  \begin{itemize}
    \item gera o espaço
      \[
        \pause
        V = \ger\left(\mathcal{B}\right)
        \pause
      \]
    \item é linearmente independente
      \[
        \pause
        c_1\vec{v}_1+\cdots+c_n\vec{v}_n = \vec{0}
        \pause
        \quad\Rightarrow\quad
        c_i = 0
      \]
  \end{itemize}
\end{frame}

%------------------------------------------------------------------------------%
\begin{frame}{Coordenadas em uma Base}
  Considerando uma base ordenada de $V$
  \[
    \pause
    \mathcal{B} = \{\vec{v}_1,\ldots,\vec{v}_n\}
  \]

  \pause
  Existe uma única representação para cada vetor \(\vec{v}\in V\)
  \[
    \pause
    \left[\,\vec{v}\,\right]_{\mathcal{B}}
    \pause =
    \begin{pmatrix}[c]
      a_{1}  \\
      \vdots \\
      a_{n}
    \end{pmatrix}
  \]
  \pause
  tal que
  \(
    \vec{v} = a_1\vec{v}_1 + \cdots + a_n\vec{v}_n
  \)
\end{frame}

%------------------------------------------------------------------------------%
\section{Bases Ortogonais}

%------------------------------------------------------------------------------%
\begin{frame}{Independência de Vetores Ortogonais}
  \onslide<+->{Vetores ortogonais
  \(
    \vec{v}_1, \ldots, \vec{v}_n
  \)}
  \onslide<+->{são \emph{linearmente independentes}}
  \begin{align*}
    \onslide<+->{\vec{0} & = c_1\vec{v}_1+\cdots+c_n\vec{v}_n \\}
    \onslide<+->{\vec{0} \cdot \vec{v}_i & = \left(c_1\vec{v}_1+\cdots+c_n\vec{v}_n\right) \cdot \vec{v}_i \\}
    \onslide<+->{0 & = c_1\vec{v}_1 \cdot \vec{v}_i + \cdots + c_i\vec{v}_i \cdot \vec{v}_i + \cdots + c_n\vec{v}_n \cdot \vec{v}_i \\}
    \onslide<+->{0 & = 0 + \cdots + c_i\vec{v}_i \cdot \vec{v}_i + \cdots + 0\\}
    \onslide<+->{0 & = c_i \norm{\vec{v}_i}^2 \\}
    \onslide<+->{c_i & = 0}
  \end{align*}
\end{frame}

%------------------------------------------------------------------------------%
\begin{frame}{Base Ortogonal}
  Uma base
  \[
    \mathcal{B} = \{\vec{v}_1, \ldots, \vec{v}_n\}
  \]
  \pause
  é \emph{ortogonal} quando
  \[
    \pause
    \vec{v}_i\cdot\vec{v}_j = 0
    \qquad
    i \neq j
  \]
\end{frame}

%------------------------------------------------------------------------------%
\begin{frame}{Coordenadas em uma Base Ortogonal}
  \onslide<+->{Se $\mathcal{B}$ é ortogonal}
  \begin{align*}
    \onslide<+->{\vec{x} & = c_1\vec{v}_1 + \cdots + c_n\vec{v}_n \\}
    \onslide<+->{\vec{x} \cdot \vec{v}_j & = \left(c_1\vec{v}_1 + \cdots + c_n\vec{v}_n\right)  \cdot \vec{v}_j \\}
    \onslide<+->{\vec{x} \cdot \vec{v}_j & = c_1\vec{v}_1 \cdot \vec{v}_j + \cdots + c_j\vec{v}_j \cdot \vec{v}_j + \cdots + c_n\vec{v}_n \cdot \vec{v}_j \\}
    \onslide<+->{\vec{x} \cdot \vec{v}_j & = 0 + \cdots + c_j\vec{v}_j \cdot \vec{v}_j + \cdots + 0\\}
    \onslide<+->{\vec{x} \cdot \vec{v}_j & = c_j \norm{\vec{v}_j}^2 \\}
    \onslide<+->{c_j & = \frac{\vec{x}\cdot\vec{v}_j}{\norm{\vec{v}_j}^{2}}}
  \end{align*}
\end{frame}

%------------------------------------------------------------------------------%
\begin{frame}{Coordenadas em uma Base Ortogonal}
  Em qualquer base ortogonal
  \[
    \vec{x} = \sum_{j=1}^{n} \;\;
    \frac{\vec{x}\cdot\vec{v}_j}{\norm{\vec{v}_j}^{2}} \;\; \vec{v}_j
  \]
\end{frame}

%------------------------------------------------------------------------------%
\section{Bases Ortonormais}

%------------------------------------------------------------------------------%
\begin{frame}{Base Ortonormal}
  Uma base
  \[
   \mathcal{B} = \{\vec{v}_1, \ldots, \vec{v}_n\}
  \]
  \pause
  é \emph{ortonormal} quando
  \[
    \pause
    \vec{v}_i\cdot\vec{v}_j = 0
    \qquad
    i\neq j
  \]
  \[
    \pause
    \norm{\vec{v}_i}=1
  \]
\end{frame}

%------------------------------------------------------------------------------%
\begin{frame}{Base Ortonormal Definição Equivalente}
  \[
    \vec{v}_i\cdot\vec{v}_j
    \pause
    = \delta_{ij}
    \pause
    =
    \begin{cases}
    1 & i = j\\
    0 & i \neq j
    \end{cases}
  \]
\end{frame}

%------------------------------------------------------------------------------%
\begin{frame}{Coordenadas em uma Base Ortonormal}
  Se $\mathcal{B}$ é ortonormal
  \[
    \pause
    \norm{\vec{v}_j}^{2} = 1
  \]

  \pause
  \[
    c_j
    = \frac{\vec{x}\cdot\vec{v}_j}{\norm{\vec{v}_j}^{2}}
    \pause
    = \vec{x} \cdot \vec{v}_j
  \]

  \pause
  \[
    \vec{x} = \sum_{j=1}^{n} \;\;
    (\vec{x}\cdot\vec{v}_j) \;\; \vec{v}_j
  \]
\end{frame}

%------------------------------------------------------------------------------%
\section{Exemplos}

%------------------------------------------------------------------------------%
\begin{question}
  Determine se
  \[
    \mathcal{B}
    = \left\{
      \vec{v}_1, \vec{v}_3, \vec{v}_3
    \right\}
    = \left\{
      \vetortri{1}{1}{0}, \,
      \vetortri{1}{-1}{0}, \,
      \vetortri{0}{0}{2}
    \right\}
  \]
  é uma base ortogonal ou ortonormal de $\R^3$
\end{question}

%------------------------------------------------------------------------------%
\begin{resolution}{Solução}
  \[
    \vec{v}_1\cdot\vec{v}_2
    \pause
    = \vetortri{1}{1}{0} \cdot \vetortri{1}{-1}{0}
    \pause
    = 1 \times 1 + 1(-1) + 0 \times 0
    \pause
    = 1 - 1 + 0
    \pause
    = 0
  \]
  \[
    \pause
    \vec{v}_1\cdot\vec{v}_3
    \pause
    = \vetortri{1}{1}{0} \cdot \vetortri{0}{0}{2}
    \pause
    = 1 \times 0 + 1 \times 0 + 0 \times 2
    \pause
    = 0
  \]
  \[
    \pause
    \vec{v}_2\cdot\vec{v}_3
    \pause
    = \vetortri{1}{-1}{0} \cdot \vetortri{0}{0}{2}
    \pause
    = 1 \times 0 + (-1)0 + 0 \times 2
    \pause
    = 0
  \]
  \pause
  Logo $\mathcal{B}$ é ortogonal
\end{resolution}

%------------------------------------------------------------------------------%
\begin{resolution}{Solução}
  \[
    \norm{\vec{v}_1}
    \pause
    = \norm{\vetortri{1}{1}{0}}
    \pause
    = \sqrt{1^2 + 1^2}
    \pause
    = \sqrt{2}
  \]
  \[
    \pause
    \norm{\vec{v}_2}
    \pause
    = \norm{\vetortri{1}{-1}{0}}
    \pause
    = \sqrt{1^2 + 1^2}
    \pause
    = \sqrt{2}
  \]
  \[
    \pause
    \norm{\vec{v}_3}
    \pause
    = \norm{\vetortri{0}{0}{2}}
    \pause
    = \sqrt{2^2}
    \pause
    = 2
  \]
  \pause
  Logo $\mathcal{B}$ não é ortonormal
\end{resolution}

%------------------------------------------------------------------------------%
\endinput

%------------------------------------------------------------------------------%
\begin{question}
  Considere a base ortogonal
  \[
    B=\{(1,1,0),(1,-1,0),(0,0,2)\}
  \]
  Determine as coordenadas de
  \[
    \vec{x}=(3,-1,4)
  \]
  na base $B$
\end{question}

%------------------------------------------------------------------------------%
\begin{resolution}{Solução}
  \onslide<+->{Se
  \(
    \vec{x} = c_1\vec{v}_1 + c_2\vec{v}_2+c_3\vec{v}_3
  \)}
\begin{columns}
\column{0.25\textwidth}
  \begin{align*}
    \onslide<+->{c_1 & = \frac{\vec{x}\cdot\vec{v}_1}{\norm{\vec{v}_1}^{2}} \\}
    \onslide<+->{    & = \frac{3-1}{2}                                      \\}
    \onslide<+->{    & = 1}
  \end{align*}

\column{0.25\textwidth}
  \begin{align*}
    \onslide<+->{c_2 & = \frac{\vec{x}\cdot\vec{v}_2} {\norm{\vec{v}_2}^{2}} \\}
    \onslide<+->{    & = \frac{3+1}{2}                                       \\}
    \onslide<+->{    & =2}
  \end{align*}
\column{0.3\textwidth}
  \begin{align*}
    \onslide<+->{c_3 & = \frac{\vec{x}\cdot\vec{v}_3}{\norm{\vec{v}_3}^{2}} \\}
    \onslide<+->{    & = \frac{8}{4}                                        \\}
    \onslide<+->{    & = 2}
  \end{align*}
\end{columns}
\end{resolution}

%------------------------------------------------------------------------------%
\begin{resolution}{Solução}
  \[
    [\vec{x}]_B=
    \begin{bmatrix}
      1 \\
      2 \\
      2
    \end{bmatrix}
  \]
\end{resolution}

%------------------------------------------------------------------------------%
\endinput


%%------------------------------------------------------------------------------%
%\section{Lista Mínima}
%
%%------------------------------------------------------------------------------%
%\begin{frame}{Lista Mínima}
%  \begin{enumerate}
%    \item
%  \end{enumerate}
%
%  \vfill
%  Atenção: A prova é baseada no livro, não nas apresentações
%\end{frame}

%------------------------------------------------------------------------------%
\end{document}
%------------------------------------------------------------------------------%
